\section{Discussion}\label{sec:discussion}
The completed temporal confirmation changes the incremental finding while preserving the family qualification. Under the frozen pooled-and-both-family criterion, the original affine envelope 20-check comparison improves pooled MAE from 1.3553 to 1.3257 ppm and strict 1 ppm exceedance from 43.06\% to 42.14\%. Fourteen of fifteen pooled B20 affine MAE changes are favorable in confirmation, compared with fifteen adverse changes in development. The stages remain separate; this reversal neither removes the adverse development evidence nor selects a new final-data winner. Historical benefit remains heterogeneous. FIS improves from affine B0 at 0.9471 to 0.9033 ppm with the primary schedule but remains worse in MAE than unchanged consensus at 0.8820 ppm. D12 and FIS D13 reversals preclude uniform-session improvement, and the envelope-schedule superiority criterion is unsupported.
\subsection{Control initialization before assigning reference value}
Historical correction already improves pooled RF performance before a new reference is bought in both stages; confirmation additionally supports an incremental aggregate benefit. The reporting identity in Eq.~\eqref{eq:decomp} exposes this distinction separately for MAE and exceedance. Retaining family results prevents a pooled improvement from concealing adverse FIS MAE. Additional-reference and schedule comparisons must use identical exposure keys and an effective same-estimator B0.

\subsection{Update information and conditional sensitivity}
Cumulative RLS can absorb new references: its gain is nonzero for informative regressors. The response depends on leverage, innovation and historical information, not merely 100 versus 20 rows. Forgetting at purchases discounts old information and ridge; it does not forget during unqueried exposures. The fixed $.95$ comparator and $.90$ sensitivity, including one native schedule, increase the scope of this development comparison without supporting a claim that any tested method is optimal. Actual state changes and leverage are archived.

In development, every pooled RF B20 affine MAE point change is adverse, while some lower-budget pooled point changes are favorable and endpoint and family patterns vary. OLS/ridge provide some small favorable incremental changes, showing dependence on upstream prediction and mask. Conditional resampling bounds help describe the frozen development paths; crossing zero does not prove absence or equivalence, and excluding zero does not confer population significance after sequential development. Quarter-path, lag-one and trend diagnostics expose temporal structure rather than assume away nonstationarity. The stated relative reference-uncertainty percentage is not an independent Gaussian noise floor for a paired contrast using common labels.

\subsection{Prediction sets and measurement assurance}
Conditional minimax midpoint behavior is distinct from empirical MAE. Projection's non-worsening property requires the actual residual to lie within the assumed interval. Hull containment cannot rescue a contradictory intersection. Frequent containment, a narrow radius and complete-array service assurance are separate properties. Confirmation records 38 false midpoint assurances among 74 flagged outputs (51.35\%) and among 3,600 total outputs (1.06\%), separately from 284 raw and 276 reported-set noncontainments. These are midpoint-envelope diagnostics, not affine uncertainty intervals. Standard conformal methods provide a useful future comparison only under their own exchangeability or sequential assumptions \cite{angelopoulos2021}; neither descriptive envelopes nor ordinary split conformal offer unconditional protection against arbitrary drift.

Traceable co-calibration requires characterized references and uncertainty propagation \cite{gruber2023,jcgm100,jcgm200}. Multiple accesses to one generated-concentration system cannot identify a shared reference bias. The current analysis is a reference-relative algorithm comparison. The fixed-D10 temporal confirmation provides distinct later-session evidence: family heterogeneity persists historically, while the original incremental aggregate criterion is met despite session reversals (Section~\ref{sec:c5}).

\section{Limitations}\label{sec:limitations}
Both stages use one chamber and array: three development assessment sessions and three separate confirmation sessions, each with twelve overlapping peer configurations. Sequential extensions were informed by earlier development results; family emphasis, forgetting, native scheduling and conditional resampling are declared extensions rather than independent preregistration. Error is relative to released nominal CO with stated relative uncertainty below 5.5\%, without exposure-level reference-error covariance or hard bounds. Certified physical uncertainty, independent-array coverage and a causal drift diagnosis are not established. Acquisition latency and cost are simulated, and resets limit the horizon of late purchases. R7/R14 and the UCI270 reserved batches remain unused; D11--D13 are evaluated only in the separately authorized fixed-D10 confirmation.

The calibrated hardware union includes units omitted from individual peer configurations; strict unseen-unit or single-sensor deployment is a different estimand. Late-core summaries omit transient and washout behavior, and anonymous positions do not establish serial/lot provenance. Inverse regression and numerical rank do not identify forward physical parameters. Empirical slopes and fixed temporal/session allowances are unvalidated future bounds. Conditional resampling is development-only; no confirmation bootstrap was performed. Ordinary block resampling requires regularity that these adaptive chronological paths may not satisfy; the saved conditional intervals neither rerun adaptive trajectories nor propagate fitting or reference uncertainty. The finite comparison does not rank all possible online methods. Complete local per-output provenance is retained; a public reproducibility deposit is not yet available.

\section{Conclusion}
The completed fixed-D10 temporal confirmation supports the prespecified incremental aggregate benefit: pooled affine B0 MAE falls from 1.3553 to 1.3257 ppm under the envelope-driven 20-check schedule, and strict 1 ppm exceedance falls from 43.06\% to 42.14\%. The same frozen descriptive criterion is met in both family aggregates, but adverse D12 and FIS D13 cells preclude uniform-session improvement and envelope-schedule superiority remains unsupported. FIS MAE improves from 0.9471 to 0.9033 ppm relative to affine B0 while remaining worse than unchanged consensus at 0.8820 ppm. Historical benefit is therefore heterogeneous, not uniformly confirmed.

The confirmation result reverses the development B20 pattern. In development, historical affine correction reduces pooled RF MAE from 0.8993 to 0.7843 ppm and strict exceedance from 34.17\% to 28.86\%, but all fifteen affine B20 MAE changes are adverse. Lower-budget and OLS/ridge sensitivities differ; development-only conditional resampling does not establish equivalence or population coverage. No confirmation bootstrap was performed. Confirmation also retains 38 false midpoint assurances among 74 flagged outputs and 38 among 3,600 total outputs, separately from interval noncontainment. These are midpoint-envelope diagnostics, not affine uncertainty intervals.

The contribution is an information-matched evaluation of established calibration rules across development and temporal confirmation. Each estimator's historical baseline, incremental contrast, family/session results, exact pairing and reference cost must remain visible. Both stages use one apparatus; the results establish neither a new calibration algorithm, certified uncertainty nor independent-array generalization.
